\documentclass{article}
\usepackage{amsmath}
\usepackage{amssymb}
\usepackage{graphicx}
\usepackage{hyperref}
\usepackage{physics}
\usepackage{tikz}
\usepackage{bm}
\usepackage{booktabs}
\usepackage{siunitx}

\title{Notes -- WIP with Corrected LICI Derivation}
\date{\today}

\begin{document}

\maketitle

\section{Introduction}

This document summarizes what we have achieved so far with our topological model for collective conical intersection points.

The basic arrangement of our system is that three identical diatomic molecules are placed in a cavity in a Cartesian coordinate system, where the axes of the molecules are parallel to the $x$, $y$, and $z$ axes. This arrangement can fundamentally create the formation of CIs in the system.

\section{Model Potentials}

Our system is described by two potential energy surfaces:
\begin{align}
V_x(x) &= x^2 \\
V_a(x) &= a \cdot (x - x_{\text{shift}})^2 + c_0
\end{align}

Where:
\begin{itemize}
    \item $a$ is a scaling parameter
    \item $x_{\text{shift}}$ is a shift parameter for the second potential
    \item $c_0$ is a constant energy offset
\end{itemize}

\section{Finding Conical Intersection Points}

A conical intersection occurs when the two potential energy surfaces meet at a point. For our model, this happens when:
\begin{equation}
V_a(x) - V_x(x) = 0
\end{equation}

\subsection{Mathematical Derivation}

Starting with the difference between the potentials:
\begin{align}
V_a(x) - V_x(x) &= a \cdot (x - x_{\text{shift}})^2 + c_0 - x^2 \\
&= (a - 1) \cdot x^2 - 2 \cdot a \cdot x_{\text{shift}} \cdot x + c_1
\end{align}

Where $c_1 = a \cdot x_{\text{shift}}^2 + c_0$.

We can rewrite this in the form of a perfect square plus a constant:
\begin{equation}
V_a(x) - V_x(x) = (a - 1) \cdot (x - x_{\text{prime}})^2 + c_1'
\end{equation}

Where $x_{\text{prime}}$ is the point where the quadratic term has its minimum:
\begin{equation}
x_{\text{prime}} = \frac{a}{a - 1} \cdot x_{\text{shift}}
\end{equation}

\section{Geometric Configuration of CI Points}

For a 3D system with coordinates $(r_1, r_2, r_3)$ representing the bond distances of three molecules, we need to find points where the potentials are equal. This leads to a specific geometric arrangement.

\subsection{Trivial and Nontrivial LICIs}

\subsubsection{Trivial LICIs}

Trivial LICIs occur when all three bond distances are equal:
\begin{equation}
(r_1, r_2, r_3) \text{ with } r_1 = r_2 = r_3 (= r_0)
\end{equation}

These form a seam with the property $r_1 = r_2 = r_3$.

\subsubsection{Nontrivial LICIs}

To find nontrivial LICIs, we search for points where two different bond distances $r_1 \neq r_2$ give equal potential differences:
\begin{equation}
(V_a - V_x)(r_1) = (V_a - V_x)(r_2)
\end{equation}

If such $r_1 \neq r_2$ exist, the following configurations provide nontrivial LICIs:
\begin{align}
\text{LICI}_1 &: (r_1, r_2, r_3) = (r_1, r_1, r_2) \\
\text{LICI}_2 &: (r_1, r_2, r_3) = (r_1, r_2, r_1) \\
\text{LICI}_3 &: (r_1, r_2, r_3) = (r_2, r_1, r_1) \\
\text{LICI}_4 &: (r_1, r_2, r_3) = (r_2, r_2, r_1) \\
\text{LICI}_5 &: (r_1, r_2, r_3) = (r_2, r_1, r_2) \\
\text{LICI}_6 &: (r_1, r_2, r_3) = (r_1, r_2, r_2)
\end{align}

\subsection{Geometric Constraints}

Since the trivial LICIs form a seam, we set up our closed paths in the plane perpendicular to this "trivial" seam. All points in this plane have the property:
\begin{equation}
r_1 + r_2 + r_3 = 3 \cdot r_0
\end{equation}
where $(r_0, r_0, r_0)$ is the trivial LICI in the plane.

It is easy to see that $\text{LICI}_1$, $\text{LICI}_2$, and $\text{LICI}_3$ lie in the same plane, and similarly, $\text{LICI}_4$, $\text{LICI}_5$, and $\text{LICI}_6$ also lie in the same plane.

\subsection{Derivation for LICI$_1$, LICI$_2$, LICI$_3$}

We now search for nontrivial LICIs of type $\text{LICI}_1$, $\text{LICI}_2$, $\text{LICI}_3$ with the additional requirement:
\begin{equation}
r_1 + r_1 + r_2 = 3 \cdot r_0
\end{equation}
for a predefined $r_0$.

From the equality condition $(V_a - V_x)(r_1) = (V_a - V_x)(r_2)$ with $r_1 \neq r_2$, we obtain:
\begin{equation}
r_1 = x_{\text{prime}} - \delta, \quad r_2 = x_{\text{prime}} + \delta
\end{equation}
where $\delta$ is an arbitrary parameter.

Applying the geometric constraint:
\begin{align}
r_1 + r_1 + r_2 &= 3 \cdot r_0 \\
3 \cdot x_{\text{prime}} - \delta &= 3 \cdot r_0 \\
\delta &= 3 \cdot (x_{\text{prime}} - r_0)
\end{align}

Substituting back to find $r_1$ and $r_2$:
\begin{align}
r_1 &= x_{\text{prime}} - \delta = 3 \cdot r_0 - 2 \cdot x_{\text{prime}} = r_0 - 2 \cdot (x_{\text{prime}} - r_0) \\
r_2 &= x_{\text{prime}} + \delta = 4 \cdot x_{\text{prime}} - 3 \cdot r_0 = r_0 + 4 \cdot (x_{\text{prime}} - r_0)
\end{align}

\subsection{Distance from Trivial LICI}

To determine how many LICIs are enclosed by a given circle, we need to know the distance between the nontrivial and trivial LICIs. The distance in the 3D coordinate space is:
\begin{align}
d_{\text{CI}} &= |r_2 - r_0| \cdot \frac{\sqrt{6}}{2} \\
&= 4 \cdot |x_{\text{prime}} - r_0| \cdot \frac{\sqrt{6}}{2} \\
&= 2 \cdot \sqrt{6} \cdot |x_{\text{prime}} - r_0|
\end{align}

\subsection{Analysis of LICI$_4$, LICI$_5$, LICI$_6$}

Similar expressions could be derived for $\text{LICI}_4$, $\text{LICI}_5$, and $\text{LICI}_6$. In this case, we would obtain:
\begin{equation}
\delta' = -3 \cdot (x_{\text{prime}} - r_0) = -\delta
\end{equation}

Taking this value and substituting into the expressions for $r_1$ and $r_2$:
\begin{align}
r_1' &= x_{\text{prime}} - \delta' = x_{\text{prime}} + \delta = r_2 \\
r_2' &= x_{\text{prime}} + \delta' = x_{\text{prime}} - \delta = r_1
\end{align}

Combining this with the definitions of the nontrivial $\text{LICI}_i$ configurations, we realize that these LICIs are actually the same as $\text{LICI}_1$, $\text{LICI}_2$, and $\text{LICI}_3$. More specifically:
\begin{align}
\text{LICI}_4 &= \text{LICI}_1 \\
\text{LICI}_5 &= \text{LICI}_2 \\
\text{LICI}_6 &= \text{LICI}_3
\end{align}

\section{Implementation in Our Code}

\subsection{Trivial LICI Implementation}

In our code, we implement the trivial LICI by setting:
\begin{align}
r_0 &= x_{\text{prime}} \\
x &= 2(x_{\text{prime}} - r_0) = 0
\end{align}

For the CI point at index $n_{\text{CI}}$, we set:
\begin{equation}
R_0[i] = 
\begin{cases}
r_0 + x + x & \text{if } i = n_{\text{CI}} \\
r_0 - x & \text{otherwise}
\end{cases}
\end{equation}

With our parameter values ($a_{V_x} = 1.0$, $a_{V_a} = 1.3$, $x_{\text{shift}} = 0.1$), this gives $x_{\text{prime}} = 0.433$ and $R_0 = [0.433, 0.433, 0.433]$ when $n_{\text{CI}} = 3$ (i.e., no specific CI point selected). This corresponds to the trivial LICI at $(r_0, r_0, r_0)$.

\subsection{Non-trivial LICI Implementation}

To find the non-trivial LICI points, we have implemented the following functions in our code:

\subsubsection{Finding Non-trivial LICIs}

The function \texttt{find\_nontrivial\_licis(aVx, aVa, x\_shift, r0=None)} calculates the three non-trivial LICI configurations using the derived formulas:

\begin{align}
x_{\text{prime}} &= \frac{a_{V_a}/a_{V_x}}{a_{V_a}/a_{V_x} - 1} \cdot x_{\text{shift}} \\
r_1 &= 3r_0 - 2x_{\text{prime}} = r_0 - 2(x_{\text{prime}} - r_0) \\
r_2 &= 4x_{\text{prime}} - 3r_0 = r_0 + 4(x_{\text{prime}} - r_0)
\end{align}

The function returns a dictionary containing:
\begin{itemize}
    \item \texttt{trivial\_lici}: $(r_0, r_0, r_0)$ - the trivial LICI
    \item \texttt{nontrivial\_licis}: List of three non-trivial LICI points
    \item \texttt{r0}, \texttt{r1}, \texttt{r2}: The calculated bond distances
    \item \texttt{x\_prime}: The position of the potential minimum
    \item \texttt{d\_CI}: Distance from trivial to non-trivial LICI
\end{itemize}

\subsubsection{Non-trivial LICI Configurations}

The three distinct non-trivial LICI configurations are:
\begin{align}
\text{LICI}_1 &: (r_1, r_1, r_2) \quad \text{(molecules 1\&2 equal, molecule 3 different)} \\
\text{LICI}_2 &: (r_1, r_2, r_1) \quad \text{(molecules 1\&3 equal, molecule 2 different)} \\
\text{LICI}_3 &: (r_2, r_1, r_1) \quad \text{(molecules 2\&3 equal, molecule 1 different)}
\end{align}

\subsubsection{Parameter Sets and Dataset Configurations}

Our code supports multiple dataset configurations for testing different scenarios:

\textbf{Dataset 2 (Trivial LICI):}

$R_0 = (r_0, r_0, r_0)$ with $r_0 = x_{\text{prime}} = 0.433$

Note: When $r_0 = x_{\text{prime}}$, the non-trivial LICI points coincide with the trivial LICI since $r_1 = r_2 = r_0$.


\textbf{Dataset 3 (Non-trivial LICI region):}

$R_0 = (r_1, r_1, r_1)$ with $r_1 = 3r_0 - 2x_{\text{prime}} = 0.433$

This dataset uses $R_0$ centered at $r_1$, which equals $r_0$ when $r_0 = x_{\text{prime}}$.


\textbf{Dataset 4 (Non-trivial LICI region):}

$R_0 = (r_2, r_2, r_2)$ with $r_2 = 4x_{\text{prime}} - 3r_0 = 0.433$

This dataset uses $R_0$ centered at $r_2$, which equals $r_0$ when $r_0 = x_{\text{prime}}$.


\subsubsection{Verification Functions}

We have implemented verification functions to ensure the correctness of the calculated LICI points:

\begin{itemize}
    \item \texttt{print\_lici\_info(lici\_data)}: Prints formatted information about the LICI points
    \item \texttt{verify\_lici\_properties(lici\_data, aVx, aVa, x\_shift, c\_const)}: Verifies that the found LICI points satisfy the CI condition $V_a(r) - V_x(r) \approx 0$ for all coordinates
\end{itemize}

The verification function checks that for each LICI configuration, the potential differences $V_a(r_i) - V_x(r_i)$ are equal for all three coordinates, confirming that they are indeed conical intersection points.

\subsection{Arrowhead Hamiltonian}

Since the cavity excites collectively, the system is described by a 4×4 arrowhead Hamiltonian matrix with the following structure:

\begin{equation}
H(\bm{R}_{\theta}) = 
\begin{pmatrix}
\hbar\omega + \sum_i V_x(\bm{R}_i) & t_{01} & t_{02} & t_{03} \\
t_{01} & V_e(\bm{R}_0) & 0 & 0 \\
t_{02} & 0 & V_e(\bm{R}_1) & 0 \\
t_{03} & 0 & 0 & V_e(\bm{R}_2)
\end{pmatrix}
\end{equation}

where:
\begin{itemize}
    \item $\hbar\omega$ is the energy quantum of the system
    \item $V_x(\bm{R}_i) = a_{Vx} |\bm{R}_i|^2$ is the potential energy function
    \item $V_e(\bm{R}_i) = \sum_j V_x(\bm{R}_j) + V_a(\bm{R}_i) - V_x(\bm{R}_i)$ is the effective potential
    \item $V_a(\bm{R}_i) = a_{Va} (|\bm{R}_i - x_{\text{shift}}|^2 + c)$ is the additional potential
    \item $t_{0i}$ are the transition dipole moments between states
\end{itemize}

The parameter vector $\bm{R}_\theta$ traces a perfect circle orthogonal to the $x=y=z$ line in 3D space, parameterized by angle $\theta$.

\subsection{Berry Connection and Berry Phase}

The Berry connection $\bm{A}_n(\bm{R})$ for the $n$-th eigenstate is defined as:

\begin{equation}
\bm{A}_n(\bm{R}) = i\langle n(\bm{R})|\nabla_{\bm{R}}|n(\bm{R})\rangle
\end{equation}

The Berry phase $\gamma_n$ is the integral of the Berry connection around a closed loop $C$ in parameter space:

\begin{equation}
\gamma_n = \oint_C \bm{A}_n(\bm{R}) \cdot d\bm{R}
\end{equation}

\subsection{Off-Diagonal Elements}

The off-diagonal elements $\tau_{nm}(\theta)$ for $n \neq m$ provide information about the coupling between different eigenstates. These elements can be used to analyze transitions between states and to identify regions in parameter space where eigenstates are strongly coupled.


\section{The Tau Matrix Method}

We call this the "tau matrix method" because the key matrix is the tau matrix in the codebase. However, this matrix is actually the matrix of the Non-adiabatic coupling terms.

In this method, we first compute $\tau$ matrix, then use it to compute the Berry phase.

\subsection{Definition of the Tau Matrix}

This method is a numerical approach for computing the Berry phase. The key element is the tau matrix, whose elements are defined as:

\begin{equation}
\tau_{nm}(\theta) = \langle \psi_m(\theta) | \frac{\partial}{\partial\theta} | \psi_n(\theta) \rangle
\end{equation}

where $|\psi_n(\theta)\rangle$ is the $n$-th eigenstate of the Hamiltonian $H(\bm{R}_\theta)$.

\section{Numerical Results and Analysis}

\subsection{Berry Phase Matrix Structure}

The numerical results show that the Berry phase matrix $\gamma_{nm}$ has a specific structure with the following features:

\begin{itemize}
    \item The diagonal elements $\gamma_{nn}$ are typically small (order $10^{-3}$ to $10^{-4}$), indicating minimal geometric phase accumulation for individual states.
    
    \item Large off-diagonal elements (order $10^0$) appear in specific positions, particularly between states 1-2, 2-1, 1-3, and 3-1.
    
    \item The values close to $\pm 2\pi$ in the [1,2] and [2,1] positions indicate complete phase winding, which is a signature of non-trivial topology.
\end{itemize}

\subsection{Physical Interpretation}

The observed Berry phase structure has important physical implications:

\begin{itemize}
    \item The system exhibits non-trivial geometric phase accumulation primarily in the subspace spanned by states 1 and 2.
    
    \item The antisymmetric nature of the off-diagonal elements ($\gamma_{12} \approx -\gamma_{21}$) is consistent with the expected behavior of a quantum system with time-reversal symmetry.
    
    \item The $2\pi$ phase difference indicates that the system's parameter space trajectory encloses a topological singularity.
\end{itemize}


\end{document}
